\chapter{レラティブ・バリューとボラティリティ}
\label{ch:rv}

\section{債券の裁定}

債券の裁定は、ほぼ同じ将来のお金の流れを持つ二つの債券（あるいは債券とデリバティブ）の価格の小さな差をとる。差は小さいので、大きなレバレッジをかける。代表的な取引は次のとおり。

\begin{itemize}
\item \textbf{発行されたばかりの銘柄と、少し前の銘柄}：最新の国債は売買が多く、少し割高になる。少し前の銘柄を買い、最新の銘柄を売り、差が縮むのを待つ。
\item \textbf{スワップ・スプレッド}：国債の利回りと、同じ年限の金利スワップの固定金利の差に賭ける。
\item \textbf{利回り曲線}：曲線の形の歪み（ある年限だけが周りより高い）に賭ける。
\item \textbf{現物と先物のベーシス}：国債の現物を買い（レポで資金を借りる）、先物を売る。
\end{itemize}

\begin{paperbox}[Duarte・Longstaff・Yu（2007）「蒸気ローラーの前の小銭か」\cite{dly2007}]
スワップ・スプレッド、利回り曲線、住宅ローン担保証券、金利のボラティリティ、資本構成（同じ会社の株と社債）の裁定の 5 つの債券の裁定の戦略を、規則で作り直して検証した。多くの戦略は有意な超過リターンを上げ、その一部は、「小さな利益を重ねて、まれに大損する」という通説とは違い、正の歪度を持っていた。分析に高度なモデルを要する（「知的な資本」の要る）戦略ほど、超過リターンが大きかった。
\end{paperbox}

\subsection{LTCM}

この種の取引で最も有名な失敗が、1998 年の LTCM である。LTCM は、債券の裁定を中心に、資本の 25 倍前後の資産を持ち（デリバティブを含めるとさらに大きい）、ロシアの債務不履行の後の「質への逃避」で、ほぼすべての取引の差が同時に広がって、数か月で資本の大半を失った。詳しくは第 \ref{ch:failures} 章で扱う。

\subsection{2020 年 3 月の国債の売り}

\begin{paperbox}[Kruttli・Monin・Petrasek・Watugala（2021）\cite{kruttli2021}]
米連邦準備制度の報告（Form PF のデータ）を使い、ヘッジファンドの米国債の取引を分析した。2020 年 3 月の新型コロナの急落で、国債の市場の変動が急に高まると、レポの資金の条件が厳しくなり、国債の取引をしていたファンドは損失を出し、国債のポジションを大きく減らした（論文は、平均的なファンドの国債のポートフォリオの損失を約 7\%、ポジションの縮小を約 2 割と報告している）。売りは、資金の調達の制約の厳しいファンドに集中した。中央銀行が国債の大量の買い入れで介入し、市場は落ち着いた。
\end{paperbox}

現物と先物のベーシスの取引は、その後も拡大し、金融の安定の観点から当局が注目している。

\section{転換社債の裁定}

転換社債（一定の条件で株に転換できる社債）を買い、その株を空売りして、株価の変動のリスクを打ち消す（デルタ・ヘッジ）。転換社債は、発行の時に割安になりやすいことが知られ、その割安さと、株価の変動から生じる利益（ガンマの利益）をとる。

Agarwal・Fung・Loon・Naik（2011）\cite{aflm2011} は、転換社債の裁定のヘッジファンドのリターンが、転換社債の市場の要因（新しい発行の量、割安さ）で説明できることを示した。新しい発行が多く割安な時期には収益が上がり、裁定の資金が集まって割安さが小さくなると、収益も小さくなる。

\subsection{遅い資本}

\begin{paperbox}[Mitchell・Pedersen・Pulvino（2007）「遅い資本」\cite{mpp2007}]
2005 年に転換社債の裁定のファンドから資金が大量に引き上げられたとき、転換社債の価格は、理論的な価値から大きく下がり、その回復には何か月もかかった。歪みを直すはずの裁定の資本は、すぐには入ってこない（遅い資本）。専門の裁定者が損をして資金を失うと、ほかの投資家が代わりに入るまでに時間がかかるからである。
\end{paperbox}

Mitchell・Pulvino（2012）\cite{mitchellpulvino2012} は、2008 年の金融危機で、プライムブローカーが担保の再利用（再担保）による資金の提供を急に止めたことが、転換社債の裁定や合併裁定のスプレッドの急拡大（裁定の崩れ）を引き起こしたことを示した。これは、裁定の収益が「資金の提供者が資金を出し続けること」に依存していることを示す。

\section{ボラティリティのリスク・プレミアム}

\subsection{オプションは平均して割高}

\begin{paperbox}[Coval・Shumway（2001）\cite{covalshumway2001}]
\SP{} の指数のオプションのリターンを調べ、プット・オプションのリターンが CAPM の予測よりずっと低いことを示した。市場のリスクを打ち消した（ベータがゼロの）アット・ザ・マネーのストラドル（コールとプットの買い）は、週に約 3\% の平均の損失を出した。変動率そのもの（あるいは急落）のリスクに、対価が支払われていることを示す。
\end{paperbox}

\begin{paperbox}[Bakshi・Kapadia（2003）\cite{bakshikapadia2003}]
オプションを買って株価の変動のリスクを打ち消したポートフォリオ（デルタ・ヘッジしたオプション）の損益が、平均して負であることを示した。市場の変動率のリスクに負の価格（変動率が上がるリスクを避けるために、投資家が対価を払う）がついていることの証拠である。
\end{paperbox}

\begin{paperbox}[Carr・Wu（2009）「分散のリスク・プレミアム」\cite{carrwu2009}]
オプションの価格から、分散スワップ（将来の実現の分散と、事前に決めた固定の分散を交換する取引）の公正な固定の値を計算し、実際に実現した分散と比べた。\SP{} などの株価指数では、実現の分散は平均してオプションが示す分散よりずっと小さく、分散スワップの買いは平均して大きな損失を出した（分散のリスク・プレミアムが大きく負）。このプレミアムは、CAPM のベータや Fama・French のファクターでは説明できない。
\end{paperbox}

\subsection{相関のリスク}

Driessen・Maenhout・Vilkov（2009）\cite{dmv2009} は、株価指数のオプションが、個別の株のオプションに比べて割高であることを示し、それが「銘柄の間の相関が上がるリスク」への対価であると解釈した。指数のオプションを売り、個別の株のオプションを買う取引（分散の取引、ディスパージョン）は、この相関のリスク・プレミアムをとる戦略である。

\subsection{ボラティリティの売りの崩れ：2018 年 2 月}

\begin{paperbox}[Augustin・Cheng・Van den Bergen（2021）「ボルマゲドン」\cite{volmageddon2021}]
2018 年 2 月 5 日、VIX（\SP{} のオプションが示す予想の変動率の指数）が 1 日で 2 倍以上（17 台から 37 台）に上がり、VIX の先物を売る形の上場商品（XIV など）は、価値の 9 割以上を失って償還された。これらの商品は、日々の値動きに応じて、引けにかけて VIX の先物を機械的に買い戻す仕組みだった。VIX が上がると、その買い戻しの量が膨らみ、VIX の先物をさらに押し上げた。著者らは、この仕組みの欠陥と、その売買が予測できたことを分析した。
\end{paperbox}

ボラティリティの売りは、保険を売る商売に似ている。平時のシャープレシオは非常に高く見えるが、それはまれな大損がまだ起きていないからである。成績を評価するには、シャープレシオではなく、最悪の損失と、その起きる確率を見る必要がある（\chapref{ch:perf}）。

\section{個人や研究で再現できるか}

\begin{itemize}
\item 債券の裁定、転換社債の裁定、ベーシス取引は、価格のデータ（債券や転換社債の日々の価格、レポの金利）、大きなレバレッジ、プライムブローカーの与信が必要で、個人や研究での再現はほぼできない。
\item ボラティリティのリスク・プレミアムの検証は、オプションの価格のデータ（有料）があれば可能である。ただし、まれな大損を含む長い期間で評価しないと、成績を過大に見る。
\item これらの戦略の研究の教訓は、収益の多くが「急落と資金の制約のリスクを引き受ける対価」であり、平時の成績からは、そのリスクが見えないことである。
\end{itemize}
